\begin{table}[htbp]
\centering
\caption{Estimated Parameters and Priors}
\label{tab:table4_priors}
\begin{tabular}{llccc}
\hline\hline
Parameter & Description & Prior & Mean & Std. Dev. \\
\hline
\multicolumn{5}{l}{\textit{Structural parameters}} \\
$n$, xi)$ & Sensitivity of spread to commodity prices & Normal & -0.199 & 0.045 \\
$x$, psi)$ & Elasticity of debt in spread & Normal & 2.800 & 0.500 \\
\\
\multicolumn{5}{l}{\textit{Commodity price process (AR(2))}} \\
$q¹_p̃$, rhoptil1)$ & Commodity price AR(2) - lag 1 & Beta & 0.800 & 0.20 \\
$-q²_p̃$, rhoptil2)$ & Commodity price AR(2) - lag 2 (magnitude) & Beta & 0.150 & 0.10 \\
$σ_p̃$, sigptil)$ & Std. dev. commodity price shock & Inv-Gamma & 0.050 & 2.00 \\
\\
\multicolumn{5}{l}{\textit{Persistence parameters ($q_i$)}} \\
$q_a$, rhoa)$ & Persistence of productivity shock (final) & Beta & 0.50 & 0.20 \\
$q_ã$, rhoatil)$ & Persistence of productivity shock (comm.) & Beta & 0.50 & 0.20 \\
$q_g$, rhog)$ & Persistence of growth shock & Beta & 0.50 & 0.20 \\
$q_s$, rhos)$ & Persistence of government spending shock & Beta & 0.50 & 0.20 \\
$q_m$, rhom)$ & Persistence of preference shock & Beta & 0.50 & 0.20 \\
$q_l$, rhol)$ & Persistence of interest rate risk shock & Beta & 0.50 & 0.20 \\
\\
\multicolumn{5}{l}{\textit{Volatility parameters ($\sigma_i$)}} \\
$σ_a$, siga)$ & Std. dev. productivity shock (final) & Inv-Gamma & 0.05 & 2 \\
$σ_ã$, sigatil)$ & Std. dev. productivity shock (comm.) & Inv-Gamma & 0.05 & 2 \\
$σ_g$, sigg)$ & Std. dev. growth shock & Inv-Gamma & 0.05 & 2 \\
$σ_s$, sigs)$ & Std. dev. government spending shock & Inv-Gamma & 0.05 & 2 \\
$σ_m$, sigm)$ & Std. dev. preference shock & Inv-Gamma & 0.05 & 2 \\
$σ_l$, sigl)$ & Std. dev. interest rate risk shock & Inv-Gamma & 0.05 & 2 \\
\hline\hline
\end{tabular}
\begin{tablenotes}
\small
\item Notes: Prior distributions for Bayesian estimation. Beta priors ensure persistence parameters are bounded in [0,1]. Inverse-Gamma priors ensure shock volatilities are strictly positive. The parameter $-q^2_{\tilde{p}}$ is estimated as a positive magnitude; the model uses the negative coefficient in the AR(2) process.
\end{tablenotes}
\end{table}
